\section{总结与延伸}

\subsection{核心要点回顾}

本文是 Anthropic Claude Code 团队基于内部数百个 Skills 的实战经验总结。核心要点包括：

\begin{enumerate}
    \item \textbf{Skills 是文件夹，不是文件}：包含脚本、资源、数据和配置，而非仅仅是一段 Markdown 指令。最有趣的 Skills 创造性地利用了文件夹结构和配置选项。
    \item \textbf{九大类别}：Library \& API Reference、Product Verification、Data Fetching \& Analysis、Business Process Automation、Code Scaffolding、Code Quality \& Review、CI/CD \& Deployment、Runbooks、Infrastructure Operations。好的 Skill 应归属单一类别。
    \item \textbf{编写原则}：聚焦非显而易见的知识、持续积累 gotchas、利用文件夹做 progressive disclosure、避免 railroading、设计好初始配置（config.json）、把 description 当触发条件。
    \item \textbf{高级特性}：通过日志/数据库实现记忆存储（使用 \texttt{CLAUDE\_PLUGIN\_DATA}）、提供辅助脚本让 Claude 专注于组合编排、通过 On Demand Hooks 在需要时激活安全检查。
    \item \textbf{分发与管理}：从仓库提交到内部插件市场的演进路径，让好的 Skills 自然涌现，注重质量审核，用 PreToolUse 钩子追踪使用数据。
\end{enumerate}

\subsection{延伸思考}

正如作者 Thariq Shihipar 所言：

\begin{quote}
Skills are incredibly powerful, flexible tools for agents, but it's still early and we're all figuring out how to use them best. Think of this more as a grab bag of useful tips that we've seen work than a definitive guide.
\end{quote}

Skills 作为 AI 智能体的扩展机制还处于早期阶段。与其把本文当作权威指南，不如视为经过实践验证的技巧合集。大多数好的 Skills 一开始只是几行文字加一个 gotcha，后来因为不断补充 Claude 遇到的新边界情况，才慢慢完善起来。

\textbf{最好的学习方式就是动手开始}（The best way to understand skills is to get started, experiment, and see what works for you）。

\subsection{拓展阅读}

\begin{itemize}
    \item Claude Code 官方文档中的 Skills 章节
    \item Skilljar 平台上关于 Agent Skills 的课程
    \item Claude Code Plugin Marketplace 文档
    \item Skill Creator --- Anthropic 发布的 Skills 创建工具
    \item Thariq Shihipar 的 X 账号（@trq212）--- 持续分享 Claude Code 使用经验
    \item 原文链接：\href{https://x.com/trq212/status/2033949937936085378}{Lessons from Building Claude Code: How We Use Skills}
\end{itemize}

%% --- Fin del contenido --- %%

\end{document}
